\begin{frontmatter}

\title{Frequent Losers as Cover Bidders in Public Procurement\texorpdfstring{\footnote{This version: March 2026.}}{}}

\author[insper]{Darcio Genicolo-Martins\corref{cor1}}
\ead{darciogm1@insper.edu.br}
\author[insper]{Paulo Furquim de Azevedo}
\affiliation[insper]{organization={INSPER},
  city={S\~ao Paulo}, country={Brazil}}
\cortext[cor1]{Corresponding author.}

\begin{abstract}
We propose frequent losers (FL)---firms that never win yet participate in
abnormally many public tenders---as a screening marker for cover bidding
in procurement auctions. Using 4.5 million tender-items from S\~ao
Paulo's BEC platform (2009--2019), we identify 2,735 FL firms via an
IQR-based threshold. FL presence is associated with 4--9\% higher
negotiated prices (OLS with item, year, and purchasing-unit fixed
effects). External validation shows 7.1\% of FL firms co-participate
with CADE-convicted cartelists, 3.5 times the rate expected by chance
($p < 0.001$). A leave-one-out instrumental variable yields a 21\%
price markup ($F = 396$), consistent with OLS attenuation. Bajari--Ye
tests reject exchangeability and conditional independence for FL bid
residuals. The screen requires only participation data, making it
deployable by any competition authority with electronic procurement
records.
\end{abstract}

\end{frontmatter}

\noindent\textbf{Keywords:} bid rigging, cover bidding, public procurement, cartel screening, frequent losers,\\ instrumental variables

\noindent\textbf{JEL No.:} D44, D73, H57, K21, L41

\bigskip
